% Historical expanded manuscript section; current evidence is in the audit notes.
\section{Thought experiments toward selection}
\label{sec:thought}

\subsection{Effective noise in a deterministic world}

The first construction, \texttt{sims/core/08}, puts each of forty coarse routes
over a block of a coupled logistic-map lattice. The lattice evolves
deterministically. A route responds to the difference between its block mean
and a slowly updated running mean, with a coupling proportional to its share.
Random numbers initialize the lattice and routes; none are drawn during their
subsequent evolution. Fixing the complete initial state fixes the trajectory.

In the retained default run, the chaotic lattice gives a final top-three flow
share of $0.62$, compared with $0.08$ for a frozen lattice. Its sampled block
means have mean lag-one autocorrelation about $0.01$, and the repeated run
reports identical summary statistics. These are demonstrations of deterministic
driving with irregular coarse fluctuations and concentrated flow, not a proof
that the coarse process is Markovian or converges to the It\^o model above.

An equal-start control removes the initial differences between coarse routes
by setting every $g_i=1$, while retaining microscopic variation in the fine
lattice. After 60,000 deterministic updates, the top-three shares are $0.615$,
$0.463$ and $0.529$ for three fine-state seeds, from $3/40=0.075$ initially.
The matched frozen controls remain at $0.075$. Fine-scale dynamics therefore
seed differentiated sharing without initial coarse-channel size differences
in these runs. The possible routes and coupling rule are specified, but the
dominant routes are not supplied as an initial pattern. This is a test of
seeding within that geometry, not a test of the full two-sided window.

The periodic control is informative: it gives an even larger top-three share,
$0.95$. The experiment therefore does not show that chaos is necessary, despite
an earlier description of that control. Low autocorrelation alone would not
establish unpredictability from all coarse information either. The contribution
of the lattice experiment is to make the absence of ongoing random input
explicit and reproducible within a flow-sharing example.

Deterministic dynamics admitting stochastic effective descriptions are already
established mathematically \cite{melbourne,gottwald}. Their role here is as a
candidate physical source for the disturbance, not a new discovery that
determinism and statistical noise can coexist. A diffusion limit for this
particular lattice, including its effective drift and stochastic convention,
remains to be derived. This matters because deterministic multiplicative
forcing need not produce a driftless It\^o limit \cite{gottwald}.

\subsection{Contact geometry and the window are separate tests}

The next construction, \texttt{sims/core/09}, recruits more lattice patches for
a route with a larger share, keeping coupling per patch fixed. In the retained
run, the variance-versus-mean-share slope is $1.01$ for sums over separate
patches. A common coherent patch gives slope $2.00$; using fixed patches with
amplitude proportional to share gives $1.99$. These are across-route diagnostic
slopes, not direct estimates of conditional infinitesimal diffusivity.

The final top-three shares are $0.17$, $0.08$ and $0.63$, respectively.
The coherent case uses the same fluctuation for all routes and so changes
cross-route correlation as well as variance scaling. It cannot isolate the
effect of moving $q$ to the upper edge. Both slopes near two are also at the
$m=1$ boundary, rather than beyond it. The lattice examples consequently
support a mechanism for effective disturbance and its scaling; they do not
demonstrate the full two-sided window in a single deterministic substrate.

The exact theorem answers a different, controlled question: for the specified
diffusion, when does pooled concentration vanish as a population fraction?
The lattice routes have a lower clipping rule and no fixed upper reflecting
cutoff, and their effective drift and shared clock have not been matched to
Eq.~\eqref{eq:model}. A next test must connect those models before treating
their results as one experiment.

\subsection{From unequal sharing to differential reproduction}

Concentration alone contains no lineage. Suppose each route also has a store
of usable work $A_i$, increasing at rate $\dot A_i=\eta J_i$, and a copy costs
$E_{\rm copy}$. With constant flow and no intervening loss, the time between
copies is $E_{\rm copy}/(\eta J_i)$. For time-varying flow, accumulated work
sets the number of completed copies, up to the remainder in the store. Thus
the ideal long-run rate is $b_i=\eta\overline J_i/E_{\rm copy}$.
This derives the rate conditional on a copying mechanism, an efficiency and a
cost; it does not derive the existence of copying machinery.

In \texttt{sims/core/18}, a funded copy replaces a randomly chosen route with
an inherited state plus a perturbation. The default run with copy cost $200$
ends at mean state $8.67\pm1.27$, compared with $4.32\pm1.17$ at cost $500$
and $1.05\pm0.19$ without copying. These are means and population standard
deviations across three runs, not confidence intervals. Fitted copying-rate
slopes are $0.0045$ and $0.0017$, against ideal values $0.0050$ and $0.0020$;
replacement can erase stored resources, and self-targeted attempts are counted.

Here the connection to selection becomes explicit: state varies, is inherited,
and affects reproduction through throughput. The simulation supplies
inheritance and replacement rules and uses random draws during evolution.
It is therefore a separate thought experiment, not a completed demonstration
of selection driven entirely by the deterministic lattice. Other repository
models investigate whether copying refreshes memory faster than restoring
dynamics erase it. Those address another necessary condition for the proposed
sequence, rather than following automatically from the concentration theorem.

\subsection{The existing heredity sequence}

The original repository distinguishes renewal of a pattern, reproduction across
lineages, and transfer through a separate carrier. A whirlpool illustrates
continuity of organisation through material turnover. Heredity of a variant asks
whether a particular deviation survives that renewal. Boundary models explore
maintained loops and reconstruction from a persistent bed; they do not establish
a general law of heredity. Channel-junction correlations also reflect spatial
continuity. A confluence viewed backwards is not itself an observed reproductive
event. Turbulent puffs that split and decay provide a physical comparison for
reproduction versus loss \cite{puffs}, without establishing inheritance of
selectable variants in the present model.

Core 12 supplies copying and tests preservation against restoration. Core 13
already implements detached tokens affecting another site's response; core 15
examines transit. These are existing toy constructions, not absent mechanisms.
Their transmission and response laws are assumptions. A small catalytic
configuration could carry a property without carrying an entire channel, and
could itself perform the catalytic function. Schr\"odinger's aperiodic crystal
is a historical image for stored arrangement \cite{schrodinger}, not a copying
mechanism. Autocatalytic RNA experiments distinguish reproduction from persistence
of composition \cite{ameta}; synthetic self-replicators selected for photocatalytic
function provide a chemical comparison \cite{photocatalysis}.

Deterministic disturbances become mutations when they alter a transmissible state
that reaches descendants. Temporary channel fluctuations need not do so. The new
checks below use explicit random draws; the deterministic lattice and inherited
carrier have not yet been joined in one construction.

\subsection{Establishment with inherited states}

The old lineage instrument counted births at a site, allowing confusion between
the founder and later occupants. It now tracks individual identity. Survival at
a finite horizon cannot establish eventual survival. Nor does a failed scalar
founder-birth count invalidate branching theory.

The new calculation uses a finite trait lattice with row transition generator
$T$, birth-attempt rates $b_i$, disappearance rates $d_i>0$, and a substochastic
offspring matrix $P$. Here $P_{ij}$ is the probability an attempt by type $i$
produces a viable daughter of type $j$; missing row mass means a failed copy.
Parents survive birth events. In a fixed background the probability $Q_i(t)$ of
extinction by time $t$, from one individual, obeys
\begin{equation}
 \dot Q=TQ+b\mathbin{\odot}Q\mathbin{\odot}(PQ-P\mathbf1)
              +d\mathbin{\odot}(\mathbf1-Q),\qquad Q(0)=0,
\end{equation}
where $\odot$ means elementwise multiplication. With perfect inheritance $P=I$,
the birth term is $b_i(Q_i^2-Q_i)$. Eventual extinction is the minimal fixed point.
This is established multitype branching mathematics \cite{branching}. For the
finite irreducible, nonsingular process used here, positive eventual survival is
equivalent to
\begin{equation}
 \rho(K)>1,\qquad K=(\operatorname{diag}d-T)^{-1}\operatorname{diag}(b)P.
\end{equation}
With $b_i=\eta J_i/E_{\rm copy}$, the corresponding copying-cost threshold is
\begin{equation}
 E_{\rm copy}<E_{\rm critical}
 =\eta\rho\!\left[(\operatorname{diag}d-T)^{-1}\operatorname{diag}(J)P\right].
\end{equation}
This applies known mathematics to measured flow, state change, loss and
transmission under the stated approximation. Exponential birth clocks are
assumed; a deterministic accumulator needs an additional work or age state.
The background is fixed while a lineage is rare, so this is not a prediction of
takeover of a finite shared flow.

Analytic one-state birth/death cases check the implementation. Another 18,000
exact-event trials check finite-time extinction, and 12,000 trials examine both
sides of the threshold. Cap hits are censored rather than counted as survival.
All 30 predictions lie within marginal 99\% Wilson intervals widened for
censoring. These are not simultaneous bounds or proof of continuum convergence.
Raw outcomes, timestep refinement and hashes are in \texttt{docs/heredity\_results/}.

\subsection{A conditional reversal of the threshold}

The retained extension takes $g\in[1,3]$, diffusion $.02g$, restoration
$-k(g-1)$, disappearance $.4$, and birth attempts $cg$. It assumes viable
transmission probability $\exp(-.02g\tau)$ over transit time $\tau$: greater
activity also damages the carrier, with a damaging event preventing a viable
copy. This is an additional hypothesis, not a consequence of the noise exponent.
Transit survival is applied at birth without an explicit travel delay.

On 64 grid intervals, restoration $k=.7$ raises the critical copying coefficient
from $.1747$ to $.3517$ at $\tau=0$, but lowers it from $1.5880$ to $1.3803$ at
$\tau=60$. The sign reversal also occurs on 16 and 32 intervals. Replacing the
damage rate $.02g$ by the constant $.02$ removes the reversal at every tested
transit time (Fig.~\ref{fig:heredity}). Restoration sacrifices flow but can improve
transmission enough to help establishment. For a fixed trait the viable birth
rate $cg\exp(-.02g\tau)$ peaks at $g=1/(.02\tau)$ if that lies in the allowed
range. The branching calculation includes changes of trait during life.

\begin{figure}[t]
\centering
\includegraphics[width=\linewidth]{heredity_threshold.pdf}
\caption{Critical copying coefficient on 64 trait intervals. Lower values make
establishment easier. Points are numerical calculations; lines guide the eye.
Activity-dependent damage allows restoration to change from a disadvantage to an
advantage as transit grows. The constant-damage control does not. These are model
predictions, not empirical observations.}
\label{fig:heredity}
\end{figure}

This is a candidate prediction of the specified couplings, distinct from the
stationary concentration window. Its originality is not established.
Reproduction--fidelity trade-offs already have experimental precedents
\cite{speedfidelity}. A physical test must measure carrier damage rather than
infer it from channel diffusivity alone.

\subsection{Matched transfer and repeated reproduction}

The new assay gives binary and continuous carriers identical initial values
$h=\pm1$ and the same destination response $g=1+.5\tanh(h)$. Destinations are
fresh: no original channel state or work store accompanies the carrier. Binary
carriers flip at rate $.005$; continuous carriers follow an Ornstein--Uhlenbeck
law with noise amplitude $.1$ and restoration $.005$ or $.3$. Exact transition
sampling avoids timestep error. At transit time 40 the destination response
difference between source types averages $.510$ for binary, $.578$ for slow
continuous, and approximately zero for fast continuous carriers (ten seeds of
10,000 packets). Permutation effects are near zero. Continuous storage can
transmit state over these finite times; universal gating necessity is not
supported. Physical storage costs are not matched.

A separate generational model funds 256 unit-cost births from total flow one
over 256 time units. Parent opportunities are proportional to flow, or equal in
a neutral control; only carrier state is transmitted and each site's response
is rebuilt. After 40 generations, positive-state fractions average $.919$ for
binary and $.966$ for slow continuous carriers, versus $.483$ and $.486$ in the
neutral control (20 seeds). Fast continuous carriers remain near $.5$. This
assumes a common funding pool and copying rules. Shuffling the recorded
parent--child pairing preserves population trajectories while removing their
association. It is a statistical null: a rising mean alone cannot identify
inheritance. Carrier origins and physical copying remain open.

\subsection{What the combined proposal commits to}

The scientific proposal is the connection between these steps, including the
bounded concentration condition. The theorem constrains this concentration
mechanism, not every possible route to Darwinian selection. It would gain substance from a construction
in which the same physical rules supply effective variation, cross both noise
boundaries, preserve inherited differences and fund reproduction. The
equal-start lattice control demonstrates a seed mechanism for differentiation
among initially equivalent possible routes. Extending that mechanism to a
material whose spatial flow paths and connectivity also develop requires a
spatial model of that response. What remains incomplete is the joining of
seeding, the controlled window test, memory and reproduction in one construction.
Nor does the prior art cited for the separate ingredients establish that
the complete proposed mechanism is already known.

The default outputs and equal-start controls are retained with source hashes
and reproduction commands in \texttt{docs/research\_examples/}. They report
finite-time toy-model behaviour. They neither establish empirical validity nor
replace the stationary checks below.
